\documentclass[tikz,border=12pt]{standalone}
\usepackage{fontspec}
\usepackage[mathrm=sym]{unicode-math}
\setmainfont{texgyrepagella-regular.otf}[
    BoldFont = texgyrepagella-bold.otf,
    ItalicFont = texgyrepagella-italic.otf,
    BoldItalicFont = texgyrepagella-bolditalic.otf
]
\setmathfont{texgyrepagella-math.otf}[AutoFakeBold = 2.5]
\usepackage{tikz}
\usetikzlibrary{arrows.meta, calc}

\begin{document}
\begin{tikzpicture}[>=Stealth, line cap=round, line join=round]

  % Parameters for potential well: V(x) = -V0 / (1 + exp((|x| - R)/a))
  \def\Vzero{3.6}
  \def\R{1.7}
  \def\a{0.22}

  % Horizontal top axis at E = 0
  \draw[line width=0.8pt] (-3.2, 0) -- (3.2, 0);

  % Vertical energy axis
  \draw[line width=0.9pt, ->] (0, -3.9) -- (0, 1.2) node[anchor=west, font=\large] {$E$};

  % Origin circle
  \filldraw[line width=0.8pt, fill=white] (0, 0) circle (2.2pt);

  % Potential well profile
  \draw[line width=1.3pt, domain=-2.8:2.8, samples=120, smooth]
    plot (\x, {-\Vzero / (1 + exp((abs(\x) - \R)/\a))});

  % Bound levels inside the well: (Energy y, half-width x)
  % E1 = -3.1
  \draw[line width=0.9pt] (-1.35, -3.1) -- (1.35, -3.1);
  % E2 = -2.6
  \draw[line width=0.9pt] (-1.55, -2.6) -- (1.55, -2.6);
  % E3 = -2.1
  \draw[line width=0.9pt] (-1.65, -2.1) -- (1.65, -2.1);
  % E4 = -1.6
  \draw[line width=0.9pt] (-1.74, -1.6) -- (1.74, -1.6);
  % E5 = Fermi level E_F = -1.1
  \draw[line width=1.0pt] (-1.85, -1.1) -- (1.85, -1.1);
  % Unoccupied levels above E_F
  \draw[line width=0.9pt] (-1.96, -0.65) -- (1.96, -0.65);
  \draw[line width=0.9pt] (-2.10, -0.25) -- (2.10, -0.25);

  % Fermi level label and dashed line to the right
  \node[anchor=east, font=\large] at (-1.9, -1.1) {$E_F$};
  \draw[line width=0.8pt, dashed] (1.85, -1.1) -- (3.3, -1.1);

  % Labels on the right side
  \node[anchor=west, align=left, font=\normalsize] at (2.2, -0.45) {empty\\[2pt]levels};
  \node[anchor=west, align=left, font=\normalsize] at (2.2, -2.0) {occupied\\[2pt]levels};

\end{tikzpicture}
\end{document}
